% Generated draft: requires mathematical and source review.
\documentclass[11pt]{article}
\usepackage[margin=25mm]{geometry}
\usepackage{fontspec}
\setmainfont{Latin Modern Roman}
\usepackage{amsmath,amssymb,mathtools,unicode-math}
\setmathfont{Latin Modern Math}
\usepackage{enumitem,needspace,xurl,hyperref}
\setlength{\parindent}{0pt}
\setlength{\parskip}{5pt}
\setlength{\emergencystretch}{3em}
\newcommand{\N}{\mathbb{N}}
\newcommand{\Z}{\mathbb{Z}}
\newcommand{\Q}{\mathbb{Q}}
\newcommand{\R}{\mathbb{R}}
\newcommand{\C}{\mathbb{C}}
\newcommand{\F}{\mathbb{F}}
\newcommand{\A}{\mathbb{A}}
\newcommand{\PP}{\mathbb{P}}
\newcommand{\E}{\mathbb{E}}
\newcommand{\eps}{\varepsilon}
\newcommand{\dd}{\,\mathrm{d}}
\begin{document}
\section*{UnsolvedMath ID: 701758}
\textbf{Source problem code:} catalogue entry
\subsection{Definitions and mathematical statement}
For a group \(G\) with identity \(e\), the order of an element \(g\in G\) is the least positive integer \(m\) such that \(g^m=e\), if such an integer exists, and is infinite otherwise. For a positive integer \(i\), write \(g^i\) for the product of \(i\) copies of \(g\).

\textbf{Conjecture (Z.-W. Sun, Kourovka Notebook Problem 20.100).} For every positive integer \(n\) and every group \(G\) containing no element of order in \(\{2,\ldots,n+1\}\), every subset \(A\subseteq G\) of cardinality \(n\) admits a bijective enumeration \(A=\{a_1,\ldots,a_n\}\) such that
\[
 a_i^{\,i}\ne a_j^{\,j}\qquad\text{for all }1\le i<j\le n.
\]
The group is not assumed to be finite or abelian, and the identity is allowed to belong to \(A\).

The supplied record states that the conjecture holds for \(n\le 3\). This small-case assertion and the record's open status are supplied source assessments, not independently verified here.
\subsection{Short English statement}
Sun conjectures that any \(n\)-element subset of a group with no elements of orders \(2,\ldots,n+1\) can be ordered so that raising its \(i\)-th element to the \(i\)-th power gives pairwise distinct elements. The supplied record reports the cases \(n\le 3\) as known.
\subsection{Sources}
\begin{itemize}
\item Local UnsolvedMath registry, permanent ID \(701758\), exact source code \texttt{catalogue entry}, title \emph{Kourovka Notebook Problem 20.100}; supplied status: open. Supplied upstream revision: \texttt{d371ee917ac9647fbb06cd6fc529deec6f3dda82}.
\item \emph{The Kourovka Notebook (Unsolved Problems in Group Theory)}, 20th issue (2022), original problem number 20.100; proposed by Z.-W. Sun in 2022. The supplied background also identifies the 21st issue edition (2026), \texttt{arXiv:1401.0300v48}.
\item The supplied collection description credits E. I. Khukhro and V. D. Mazurov as editors of the Kourovka Notebook collection covering Issues 1--21 (1965--2026). These bibliographic details and the reported small-case result are supplied attribution and have not been independently checked.
\end{itemize}
\end{document}
